% BEGIN PROTES BENCHMARK ADDITION
\paragraph{Two PROTES access models.}
We additionally compare with PROTES \cite{batsheva2023protes}, separating
access to an exact feasible-set indicator from initialization using only
feasible observations. These experiments concern binary $n\times n$
assignment, with $N=n^2$ variables in row-major order. They are distinct from
the preceding trained-kernel comparison with ConstrainTNet.
All runs use float64 and one logical CPU on the same local machine;
\cref{app:protes_protocol} specifies the construction, timing, training,
software versions, and recorded sample and iteration counts.

\paragraph{Exact full indicator: storage and sampling.}
When all feasible charge states are known, both representations describe
the complete set of $n!$ assignments. Their minimal total bond dimensions
coincide; the advantage of the symmetric representation is its block sparsity,
not a smaller total exact TT rank. In \cref{fig:protes_exact}, dense storage
grows substantially faster than storage of the permitted scalar blocks.
Actual Julia object sizes include substantial metadata: at $n=16$, the full
$U(1)$ object occupies $1.59$ GiB, compared with $18.0$ MiB of numerical blocks.
The corresponding dense variable-rank cores would require approximately
$472$ GiB; this last number is a storage calculation, not an allocation.

For sampling we construct the complete indicator in both libraries and
use the same uniform distribution over assignments, with no training.
All observed samples are feasible. At $n=8$ and $n=10$, OhMyU1 generates
approximately $2.1\cdot10^4$ and $7.8\cdot10^3$ samples per second,
respectively, versus $2.3\cdot10^3$ and $1.1\cdot10^2$ for the more
memory-efficient \texttt{protes\_general}. The resulting speedups are
approximately $9\times$ and $71\times$; small instances do not exhibit
universal superiority. These measurements compare particular dense and
block-sparse implementations, rather than establishing a lower bound on
all possible PROTES implementations.

\begin{figure*}[t]
\centering
\begin{subfigure}[t]{0.48\textwidth}
\centering\includegraphics[width=\linewidth]{pictures/benchmark_protes_memory_julia.pdf}
\caption{Storage of the full assignment indicator.}
\end{subfigure}\hfill
\begin{subfigure}[t]{0.48\textwidth}
\centering\includegraphics[width=\linewidth]{pictures/benchmark_protes_sampling_julia.pdf}
\caption{Sampling from the full uniform feasible distribution.}
\end{subfigure}
\caption{Exact-indicator benchmarks, with no training or finite-data initialization.
(a) The two PROTES curves count float64 coefficients in dense TT cores with
uniform or minimal variable ranks; the scalar-block $U(1)$ curve counts only
allowed coefficients. Markers measure the complete Julia object, including
charges and dictionaries, through $n=16$. These are storage sizes, not peak
training memory. (b) Median sampling throughput over three seeds; shading
shows the observed minimum--maximum range. Each trial accumulates at least
one second of timed generation in batches of 100. Counts include repeated
draws. Construction, canonicalization, compilation, and diagnostic bookkeeping
are excluded. Complete dense models are actually constructed only for
$4\le n\le10$; complete $U(1)$ models are constructed through $n=16$.
The distributions and binary variable order are identical.}
\label{fig:protes_exact}
\end{figure*}

\paragraph{Initialization from feasible observations.}
We next provide the same 400 feasible observations and objective instance
to PROTES and $U(1)$-EELS, without providing an indicator or permutation
encoding to PROTES. OhMyU1 retains access to $A,b$, an intentional difference
in structural information. EELS uses $\gamma=1$, sector degeneracy one,
one training sweep at learning rate $0.05$, and up to 10000 generated points
per outer iteration. All methods run until a common time budget, including
initial fitting; the number of iterations is not fixed.
Five prespecified PROTES configurations vary rank, learning rate, batch size,
Adam steps, and initial-data passes (\cref{tab:protes_settings}).

\Cref{fig:protes_data_short,fig:protes_data_long} count distinct feasible
points absent from the initial dataset, rather than repeated draws.
Longer initial fitting improves PROTES: at $n=8$, \texttt{fit20} finds
12--25 new points for linear objectives and 17--25 for quadratic objectives
in 50 seconds, while EELS finds 5157--6697 and 5240--7043, respectively.
Thus, inability to generate any feasible point is not a robust conclusion;
the observed limitation is rare feasible generation under this data-based
initialization and update rule. New points need not improve the objective.

\begin{figure*}[t]
\centering\includegraphics[width=\textwidth]{pictures/benchmark_protes_u1_5s_julia.pdf}
\caption{New distinct feasible assignments within a 5-second budget, for
linear (left) and quadratic (right) objectives, versus assignment size $n$.
Each point is one matched diagnostic instance/seed; all methods receive the
same 400 initial observations. PROTES settings are in
\cref{tab:protes_settings}; $U(1)$-EELS uses the fixed settings above.
Initialization, fitting, sampling, and scoring count toward the budget.
At $n=4$ the initial data already cover all 24 solutions, so zero novelty is
unavoidable. Both panels share the ordinate transformation $\log_{10}(1+z)$,
with ticks labelled by the original count $z$, allowing zeros to be shown.
No uncertainty estimate can be obtained from this single-seed screen.}
\label{fig:protes_data_short}
\end{figure*}

\begin{figure*}[t]
\centering\includegraphics[width=\textwidth]{pictures/benchmark_protes_u1_50s_julia.pdf}
\caption{The same data-initialized comparison at $n=8$ ($N=64$), with a
50-second budget and three matched instance/seed pairs for each objective.
Dots are individual counts of new distinct feasible points; orange bars
and red diamonds show the arithmetic mean and median of the original counts.
Both panels use the same $\log_{10}(1+z)$ ordinate and limits. All five PROTES
configurations are reported, not only the best observed configuration.
These are trained, time-limited runs, unlike the untrained full-indicator
sampling experiment in \cref{fig:protes_exact}.}
\label{fig:protes_data_long}
\end{figure*}

\paragraph{Scope of the comparison.}
The fraction $n!/2^{n^2}$ of feasible binary strings decreases rapidly,
making sparsity a plausible explanation for the difficulty of learning
feasibility from observations; this experiment does not isolate that cause
from the initialization and learning rule. It does not refute constrained
PROTES supplied with an exact indicator, or demonstrate an advantage under
identical access to constraint information. The findings are specific to
this assignment encoding, parameter grid, budgets, and local implementations.
They do not extend automatically to a small number of constraints, simple
constraints with low-rank indicators, different encodings or orderings,
sparse adaptations of PROTES, or other hardware. Nor do these diagnostic
instances constitute a new held-out SOTA evaluation or a comparison of
final optimization quality.
\clearpage
% END PROTES BENCHMARK ADDITION
